% =====================================================================
% IBFE Tutorials: from conservation laws to Navier--Stokes
% =====================================================================
\documentclass[aspectratio=169,11pt]{beamer}

% Match the visual language of the SIAM Waves talk.
\usetheme{Boadilla}
\usecolortheme{seahorse}
\useinnertheme{rectangles}
\useoutertheme{infolines}

\setbeamercolor{structure}{fg=blue!55!black}
\setbeamercolor{title}{fg=white,bg=blue!55!black}
\setbeamercolor{frametitle}{fg=blue!55!black}
\setbeamercolor{block title}{bg=blue!12,fg=blue!55!black}
\setbeamercolor{block body}{bg=blue!4}
\setbeamertemplate{navigation symbols}{}
\setbeamerfont{frametitle}{series=\bfseries,size=\large}
\setbeamerfont{title}{series=\bfseries,size=\Large}

\usepackage[T1]{fontenc}
\usepackage[utf8]{inputenc}
\DeclareFontShape{T1}{cmss}{b}{n}{<->ssub*cmss/bx/n}{}
\usepackage{graphicx}
\usepackage{amsmath,amssymb,bm,mathtools}
\usepackage{booktabs}
\usepackage{xcolor}
\usepackage{tcolorbox}
\tcbuselibrary{skins}
\usepackage{hyperref}

\graphicspath{{../pictures/0-1/}{pictures/0-1/}}
\hypersetup{colorlinks=true,urlcolor=blue!55!black}

\newcommand{\hi}[1]{\textcolor{blue!55!black}{\textbf{#1}}}
\newcommand{\warn}[1]{\textcolor{red!70!black}{\textbf{#1}}}
\newcommand{\soft}[1]{\textcolor{gray!50!black}{#1}}
\newcommand{\uvec}{\bm{u}}
\newcommand{\xvec}{\bm{x}}
\newcommand{\nvec}{\bm{n}}
\newcommand{\bvec}{\bm{b}}
\newcommand{\tvec}{\bm{t}}
\newcommand{\stress}{\bm{\sigma}}
\newcommand{\taustress}{\bm{\tau}}
\newcommand{\Dtensor}{\mathbf{D}}
\newcommand{\Wtensor}{\mathbf{W}}
\newcommand{\Identity}{\mathbf{I}}

\newtcolorbox{punchbox}{
  enhanced,
  colback=blue!4,
  colframe=blue!55!black,
  boxrule=0.6pt,
  arc=2pt,
  left=4pt,right=4pt,top=3pt,bottom=3pt
}

\title[Navier--Stokes]{From conservation laws to Navier--Stokes}
\subtitle{Kinematics, stress, momentum, and boundary data}
\author[Miller]{Laura A. Miller}
\institute[UArizona]{Department of Mathematics, University of Arizona}
\date{}

\begin{document}

% ----- 1. Title -----
\begin{frame}[plain]
  \begin{beamercolorbox}[sep=10pt,center]{title}
    {\bfseries\Large From conservation laws to Navier--Stokes}\\[3pt]
    {\normalsize Kinematics, stress, momentum, and boundary data}
  \end{beamercolorbox}
  \vspace{1.4em}
  \begin{columns}[T,onlytextwidth]
  \begin{column}{0.54\textwidth}
    \textbf{\large Laura A. Miller}\\[4pt]
    Department of Mathematics\\
    University of Arizona\\[14pt]
    {\small IBFE Tutorials}
  \end{column}
  \begin{column}{0.46\textwidth}
    \begin{punchbox}
    \centering
    \[
      \rho\left(
      \frac{\partial\uvec}{\partial t}
      +(\uvec\!\cdot\!\nabla)\uvec
      \right)
      =-\nabla p+\mu\nabla^2\uvec+\rho\bvec
    \]
    \[
      \nabla\!\cdot\uvec=0
    \]
    \end{punchbox}
  \end{column}
  \end{columns}
  \vfill
  \centering{\tiny\color{black!55}AI tools were used to assist with formatting, organization, and clarification of the instructor's original notes.}
\end{frame}

% ----- Original course-resource sequence -----
\begin{frame}{Gentle introductions to fluid dynamics}
\scriptsize
\begin{itemize}\itemsep7pt
  \item \textbf{MIT OpenCourseWare: Introduction to Fluid Mechanics (1.060)}\\
  Conceptual development and physical intuition before detailed analysis.\\
  \url{https://ocw.mit.edu/courses/1-060-engineering-mechanics-ii-fluid-dynamics-fall-2013/}
  \item \textbf{Harvard Natural Sciences Lectures: fluids and biological flows}\\
  Demonstrations with biological context.\\
  \url{https://projects.iq.harvard.edu/sciencelectures}
  \item \textbf{NPTEL: Fluid Mechanics for Biological Systems}\\
  A course aimed at life-science audiences.\\
  \url{https://nptel.ac.in/courses/102103084}
\end{itemize}
\end{frame}

\begin{frame}{Advanced university courses}
\scriptsize
\begin{itemize}\itemsep6pt
  \item \textbf{MIT OpenCourseWare: 2.25 Advanced Fluid Mechanics}\\
  \url{https://ocw.mit.edu/courses/2-25-advanced-fluid-mechanics-fall-2013/}
  \item \textbf{MIT OpenCourseWare: 2.06 Fluid Dynamics}\\
  \url{https://ocw.mit.edu/courses/2-06-fluid-dynamics-spring-2020/}
  \item \textbf{Cambridge University Engineering Department: Fluid Mechanics I and II notes}\\
  \url{https://www-civ.eng.cam.ac.uk/cued-lecture-notes/fluids/}
  \item \textbf{Caltech GALCIT: fluid-mechanics education resources}\\
  \url{https://www.galcit.caltech.edu/education}
\end{itemize}
\end{frame}

\begin{frame}{Video lectures and flow visualization}
\scriptsize
\begin{itemize}\itemsep9pt
  \item \textbf{Stanford computational fluid dynamics lectures}\\
  \url{https://www.youtube.com/playlist?list=PLd3hlSJsX_ImKP68wfKjvRlKq2Ot3PuuX}
  \item \textbf{NPTEL: IIT Kanpur Fluid Mechanics}\\
  \url{https://nptel.ac.in/courses/112104118}
  \item \textbf{University of Colorado Boulder: Flow Visualization}\\
  \url{https://www.flowvis.org/}
\end{itemize}
\end{frame}

\begin{frame}{Additional references and demonstrations}
\scriptsize
\begin{itemize}\itemsep9pt
  \item \textbf{L. C. Evans: partial differential equations notes}\\
  \url{https://math.ucdavis.edu/~evans/}
  \item \textbf{Dabiri Laboratory: bio-inspired fluid dynamics videos}\\
  \url{https://dabirilab.com/videos/}
  \item \textbf{APS Division of Fluid Dynamics: tutorial videos}\\
  \url{https://www.aps.org/units/dfd/meetings/video.cfm}
\end{itemize}
\end{frame}

% ----- 2. Roadmap -----
\begin{frame}{Roadmap}
\large
\begin{enumerate}\itemsep8pt
  \item \textbf{Kinematics:} connect fixed-point measurements to the acceleration of a moving fluid parcel.
  \item \textbf{Conservation:} derive continuity from mass conservation.
  \item \textbf{Forces:} represent pressure and viscous surface forces with the Cauchy stress tensor.
  \item \textbf{Constitutive law:} specialize the momentum balance to an incompressible Newtonian fluid.
  \item \textbf{Mathematical problem:} pair the equations with compatible initial and boundary data.
\end{enumerate}
\end{frame}

% ----- Basic notation retained from the original tutorial -----
\begin{frame}{Basic variables and units}
\begin{columns}[T,onlytextwidth]
\begin{column}{0.49\textwidth}
\textbf{Position and velocity}
\[
  \xvec=(x_1,x_2,x_3)=(x,y,z),
\]
\[
  \uvec=(u_1,u_2,u_3)=(u,v,w).
\]
\end{column}
\begin{column}{0.51\textwidth}
\begin{center}
\renewcommand{\arraystretch}{1.25}
\begin{tabular}{lll}
\toprule
Quantity & Symbol & SI unit \\
\midrule
density & $\rho$ & $\mathrm{kg\,m^{-3}}$ \\
pressure & $p$ & $\mathrm{Pa}$ \\
dynamic viscosity & $\mu$ & $\mathrm{Pa\,s}$ \\
kinematic viscosity & $\nu=\mu/\rho$ & $\mathrm{m^2\,s^{-1}}$ \\
body acceleration & $\bvec$ & $\mathrm{m\,s^{-2}}$ \\
\bottomrule
\end{tabular}
\end{center}
\end{column}
\end{columns}
\begin{punchbox}
\small A body-force density has units of force per volume. Writing it as $\rho\bvec$ distinguishes that density from the acceleration $\bvec$.
\end{punchbox}
\end{frame}

\begin{frame}{Cartesian vector differential operators}
\small
\begin{columns}[T,onlytextwidth]
\begin{column}{0.50\textwidth}
\textbf{Gradient of a scalar}
\[
  \nabla\phi=
  \left(\frac{\partial\phi}{\partial x},
  \frac{\partial\phi}{\partial y},
  \frac{\partial\phi}{\partial z}\right).
\]
\textbf{Divergence of a vector}
\[
  \nabla\cdot\uvec=
  \frac{\partial u}{\partial x}
  +\frac{\partial v}{\partial y}
  +\frac{\partial w}{\partial z}.
\]
\end{column}
\begin{column}{0.50\textwidth}
\textbf{Curl of a vector}
\[
  \nabla\times\uvec=
  \left(
  \frac{\partial w}{\partial y}-\frac{\partial v}{\partial z},
  \frac{\partial u}{\partial z}-\frac{\partial w}{\partial x},
  \frac{\partial v}{\partial x}-\frac{\partial u}{\partial y}
  \right).
\]
\textbf{Vector Laplacian}
\[
  \nabla^2\uvec=(\nabla^2u,\nabla^2v,\nabla^2w).
\]
\end{column}
\end{columns}
\end{frame}

\begin{frame}{Summary of notation}
\begin{center}
\renewcommand{\arraystretch}{1.28}
\begin{tabular}{ll}
\toprule
Symbol & Meaning \\
\midrule
$\xvec=(x,y,z)$ & spatial position \\
$\uvec=(u,v,w)$ & velocity field \\
$p$, $\rho$, $\mu$ & pressure, density, dynamic viscosity \\
$\nabla\phi$ & gradient of a scalar \\
$\nabla\cdot\uvec$ & divergence or local volume-expansion rate \\
$\nabla\times\uvec$ & vorticity vector \\
$\nabla^2\uvec$ & componentwise Laplacian \\
$D/Dt$ & rate following a fluid parcel \\
\bottomrule
\end{tabular}
\end{center}
\end{frame}

% ----- 3. Descriptions -----
\begin{frame}{Eulerian and material descriptions}
\begin{columns}[T,onlytextwidth]
\begin{column}{0.50\textwidth}
\textbf{Eulerian description}

The field $\uvec(\xvec,t)$ gives the velocity observed at each fixed spatial point $\xvec$.

\medskip
This is the natural description for fluid equations on a stationary grid.
\end{column}
\begin{column}{0.50\textwidth}
\textbf{Material description}

A parcel labeled by $\bm{\xi}$ follows a trajectory $\mathbf{X}(\bm{\xi},t)$ satisfying
\[
  \frac{d\mathbf{X}}{dt}
  =\uvec(\mathbf{X}(\bm{\xi},t),t).
\]
This is the natural description for following matter.
\end{column}
\end{columns}
\begin{punchbox}
\centering
\small The material derivative translates an Eulerian field into the rate experienced along a parcel trajectory.
\end{punchbox}
\end{frame}

% ----- 4. Scalar material derivative -----
\begin{frame}{Material derivative of a scalar}
Let $\phi(\xvec,t)$ be temperature, concentration, or another scalar field. Along a parcel path,
\begin{align*}
  \frac{d}{dt}\phi(\mathbf{X}(t),t)
  &=\frac{\partial\phi}{\partial t}
    +\frac{d\mathbf{X}}{dt}\cdot\nabla\phi \\
  &=\frac{\partial\phi}{\partial t}
    +\uvec\cdot\nabla\phi.
\end{align*}
Therefore
\[
  \boxed{\frac{D\phi}{Dt}
  =\frac{\partial\phi}{\partial t}
  +\uvec\cdot\nabla\phi}.
\]
\begin{columns}[T,onlytextwidth]
\begin{column}{0.50\textwidth}
\hi{Local change}: $\partial\phi/\partial t$ at a fixed point.
\end{column}
\begin{column}{0.50\textwidth}
\hi{Advective change}: motion through spatial gradients.
\end{column}
\end{columns}
\end{frame}

\begin{frame}{Material derivative in Cartesian coordinates}
For a scalar field $\phi(x,y,z,t)$ and parcel velocity $\uvec=(u,v,w)$,
\[
  \frac{d\phi}{dt}
  =\frac{\partial\phi}{\partial t}
  +\frac{dx}{dt}\frac{\partial\phi}{\partial x}
  +\frac{dy}{dt}\frac{\partial\phi}{\partial y}
  +\frac{dz}{dt}\frac{\partial\phi}{\partial z}.
\]
Substituting $dx/dt=u$, $dy/dt=v$, and $dz/dt=w$ gives
\[
  \boxed{
  \frac{D\phi}{Dt}
  =\frac{\partial\phi}{\partial t}
  +u\frac{\partial\phi}{\partial x}
  +v\frac{\partial\phi}{\partial y}
  +w\frac{\partial\phi}{\partial z}}
\]
This is the component expansion of
$D\phi/Dt=\partial_t\phi+\uvec\cdot\nabla\phi$.
\end{frame}

% ----- 5. Acceleration -----
\begin{frame}{Acceleration of a fluid parcel}
Apply the material derivative componentwise to velocity:
\[
  \boxed{
  \bm{a}
  =\frac{D\uvec}{Dt}
  =\frac{\partial\uvec}{\partial t}
  +(\uvec\cdot\nabla)\uvec
  }.
\]
In index notation,
\[
  a_i=\frac{\partial u_i}{\partial t}
  +u_j\frac{\partial u_i}{\partial x_j},
\]
where repeated indices are summed.

\begin{punchbox}
\small The advective term is nonlinear because the velocity field transports itself. A steady flow can still accelerate parcels when velocity varies in space.
\end{punchbox}
\end{frame}

\begin{frame}{Vector material derivative in components}
For $\uvec=(u,v,w)$, apply $D/Dt$ to each component:
\begin{align*}
  \frac{Du}{Dt}
  &=\frac{\partial u}{\partial t}
    +u\frac{\partial u}{\partial x}
    +v\frac{\partial u}{\partial y}
    +w\frac{\partial u}{\partial z},\\[4pt]
  \frac{Dv}{Dt}
  &=\frac{\partial v}{\partial t}
    +u\frac{\partial v}{\partial x}
    +v\frac{\partial v}{\partial y}
    +w\frac{\partial v}{\partial z},\\[4pt]
  \frac{Dw}{Dt}
  &=\frac{\partial w}{\partial t}
    +u\frac{\partial w}{\partial x}
    +v\frac{\partial w}{\partial y}
    +w\frac{\partial w}{\partial z}.
\end{align*}
\[
  \frac{D\uvec}{Dt}=
  \left(\frac{Du}{Dt},\frac{Dv}{Dt},\frac{Dw}{Dt}\right).
\]
\end{frame}

\begin{frame}{Local and advective acceleration}
\[
  \frac{D\uvec}{Dt}
  =\underbrace{\frac{\partial\uvec}{\partial t}}_{\text{local acceleration}}
  +\underbrace{(\uvec\cdot\nabla)\uvec}_{\text{advective acceleration}}.
\]
\begin{columns}[T,onlytextwidth]
\begin{column}{0.50\textwidth}
The local term measures change at a fixed spatial point.
\end{column}
\begin{column}{0.50\textwidth}
The advective term measures change as a parcel moves through a nonuniform velocity field.
\end{column}
\end{columns}
\begin{punchbox}
\small A flow can be steady, $\partial_t\uvec=0$, while individual parcels accelerate. Flow through a narrowing passage is a familiar example.
\end{punchbox}
\end{frame}

% ----- Incompressibility: geometric meaning -----
\begin{frame}{Physical meaning of incompressibility}
Imagine a small dyed blob that contains the same fluid parcels as it moves.

\medskip
The blob may translate, rotate, stretch, and shear. For an incompressible material, however, its volume remains constant:
\[
  \boxed{\frac{d}{dt}|\Omega(t)|=0}.
\]
\begin{punchbox}
\small Incompressibility restricts local volume change. It does not require the blob to retain its shape, and it does not mean that the velocity is spatially uniform.
\end{punchbox}
The next step expresses the blob's volume change in terms of flux through its moving boundary.
\end{frame}

\begin{frame}{Volume flux and the divergence theorem}
Let $\Omega(t)$ be a material region with outward unit normal $\nvec$. Its boundary moves with the fluid, so
\[
  \frac{d}{dt}|\Omega(t)|
  =\oint_{\partial\Omega(t)}\uvec\cdot\nvec\,dS.
\]
The divergence theorem converts this surface flux into a volume integral:
\[
  \oint_{\partial\Omega(t)}\uvec\cdot\nvec\,dS
  =\int_{\Omega(t)}\nabla\cdot\uvec\,dV.
\]
If every material region preserves volume, then
\[
  \boxed{\nabla\cdot\uvec=0}.
\]
\begin{punchbox}
\small Positive divergence means local expansion; negative divergence means local compression.
\end{punchbox}
\end{frame}

% ----- 6. Mass integral -----
\begin{frame}{Mass conservation for a material region}
Let $\Omega(t)$ contain the same fluid parcels for all time. Its mass is
\[
  M(t)=\int_{\Omega(t)}\rho(\xvec,t)\,dV.
\]
Mass conservation requires
\[
  \frac{d}{dt}\int_{\Omega(t)}\rho\,dV=0.
\]
The Reynolds transport theorem gives
\[
  \frac{d}{dt}\int_{\Omega(t)}\rho\,dV
  =\int_{\Omega(t)}
  \left[
  \frac{\partial\rho}{\partial t}
  +\nabla\cdot(\rho\uvec)
  \right]dV.
\]
Because the material region is arbitrary, the integrand must vanish.
\end{frame}

% ----- 7. Continuity -----
\begin{frame}{Continuity and incompressibility}
The differential mass balance is
\[
  \boxed{
  \frac{\partial\rho}{\partial t}
  +\nabla\cdot(\rho\uvec)=0
  }
  \qquad\Longleftrightarrow\qquad
  \boxed{
  \frac{D\rho}{Dt}+\rho\nabla\cdot\uvec=0
  }.
\]

For an incompressible material, parcel density does not change:
\[
  \frac{D\rho}{Dt}=0
  \qquad\Longrightarrow\qquad
  \boxed{\nabla\cdot\uvec=0}.
\]

The volume of a material region satisfies
\[
  \frac{d}{dt}|\Omega(t)|
  =\oint_{\partial\Omega(t)}\uvec\cdot\nvec\,dS
  =\int_{\Omega(t)}\nabla\cdot\uvec\,dV.
\]

\soft{\footnotesize Uniform constant density is a stronger assumption than incompressibility: a divergence-free flow may advect spatial density variations.}
\end{frame}

% ----- Pressure forces and the inviscid model -----
\begin{frame}{Pressure traction on a fluid parcel}
Let $\Omega(t)$ be a small material region with outward unit normal $\nvec$.

\medskip
Pressure is an isotropic compressive stress: its magnitude is independent of the orientation of the surface. The traction exerted by the surrounding fluid on the parcel is
\[
  \boxed{\tvec_p(\nvec)=-p\nvec}.
\]
Thus the force on a surface element $dS$ is $-p\nvec\,dS$.

\begin{punchbox}
\small Because $\nvec$ points out of the parcel, the minus sign makes a positive pressure act inward.
\end{punchbox}
\end{frame}

\begin{frame}{Net pressure force}
Adding the pressure traction over the parcel boundary gives
\[
  \boxed{
  \mathbf{F}_p=-\oint_{\partial\Omega(t)}p\nvec\,dS}.
\]
For constant pressure, contributions on opposite sides cancel and the net force is zero.

\medskip
For spatially varying pressure, the higher-pressure side pushes more strongly. The resultant points toward lower pressure.

\begin{punchbox}
\small The opposite force, exerted by the parcel on its surroundings across the same interface, follows from Newton's third law.
\end{punchbox}
\end{frame}

\begin{frame}{From surface pressure to a volume force}
Apply the divergence theorem componentwise. For each Cartesian basis vector $\bm e_i$,
\[
  \oint_{\partial\Omega}(p\bm e_i)\cdot\nvec\,dS
  =\int_\Omega\nabla\cdot(p\bm e_i)\,dV
  =\int_\Omega\frac{\partial p}{\partial x_i}\,dV.
\]
Stacking the components gives
\[
  \oint_{\partial\Omega}p\nvec\,dS
  =\int_\Omega\nabla p\,dV.
\]
Therefore
\[
  \boxed{\mathbf{F}_p=-\int_\Omega\nabla p\,dV}.
\]
\end{frame}

\begin{frame}{Small-parcel limit and sign check}
If $\nabla p$ is nearly constant within a parcel of volume $\delta V$,
\[
  \boxed{\mathbf{F}_{p,\mathrm{parcel}}
  \approx-(\nabla p)\,\delta V}.
\]
The force per unit volume is therefore $-\nabla p$.

\medskip
For the same interface and orientation, Newton's third law gives the opposite resultant exerted by the parcel on the surrounding fluid:
\[
  \mathbf{F}_{p,\mathrm{surroundings}}
  \approx+(\nabla p)\,\delta V.
\]
\begin{punchbox}
\small Always state whose force is being calculated and which region defines the outward normal.
\end{punchbox}
\end{frame}

\begin{frame}{Ideal-fluid specialization used here}
The introductory Euler model combines three assumptions:
\begin{itemize}\itemsep7pt
  \item \textbf{inviscid stress}: $\stress=-p\Identity$;
  \item \textbf{incompressible motion}: $\nabla\cdot\uvec=0$;
  \item \textbf{uniform constant density}: $\rho$ is the same for every parcel and time.
\end{itemize}
Pressure supplies the surface traction
\[
  \tvec_p=-p\nvec.
\]
\begin{punchbox}
\small Inviscid, incompressible, and constant-density are distinct assumptions. Compressible Euler equations retain the inviscid stress model while allowing density to vary.
\end{punchbox}
\end{frame}

\begin{frame}{Momentum of a dyed parcel: applied forces}
For a small material parcel of volume $\delta V$,
\[
  \mathbf{F}_p
  =-\oint_{\partial\Omega}p\nvec\,dS
  =-\int_{\delta V}\nabla p\,dV
  \approx-(\nabla p)\,\delta V.
\]
If gravity is the body acceleration, then
\[
  \mathbf{F}_g
  =\int_{\delta V}\rho\bm g\,dV
  \approx\rho\bm g\,\delta V.
\]
Hence
\[
  \boxed{\mathbf{F}_{\mathrm{total}}
  \approx(-\nabla p+\rho\bm g)\,\delta V}.
\]
\end{frame}

\begin{frame}{Momentum of a dyed parcel: Newton's law}
The parcel mass is $\rho\,\delta V$, and its acceleration is the material derivative $D\uvec/Dt$. Newton's second law gives
\[
  \rho\,\delta V\,\frac{D\uvec}{Dt}
  =(-\nabla p+\rho\bm g)\,\delta V.
\]
Canceling $\delta V$ yields
\[
  \boxed{
  \rho\frac{D\uvec}{Dt}=-\nabla p+\rho\bm g}.
\]
\begin{punchbox}
\small The same result follows exactly from an integral momentum balance for an arbitrary material region; the small-parcel argument supplies the physical picture.
\end{punchbox}
\end{frame}

\begin{frame}{Euler equations}
Using $D/Dt=\partial_t+\uvec\cdot\nabla$, the inviscid momentum balance is
\[
  \boxed{
  \frac{\partial\uvec}{\partial t}
  +(\uvec\cdot\nabla)\uvec
  =-\frac{1}{\rho}\nabla p+\bvec}.
\]
From here the general body acceleration $\bvec$ replaces the gravitational acceleration used in the parcel derivation (with $\bvec=\bm g$ for gravity).

For the constant-density incompressible specialization,
\[
  \boxed{\nabla\cdot\uvec=0}.
\]
\begin{punchbox}
\small Pressure adjusts to enforce the divergence constraint while the momentum equation determines parcel acceleration.
\end{punchbox}
\end{frame}

\begin{frame}{Euler equations in Cartesian components}
\footnotesize
For $\uvec=(u,v,w)$ and $\bvec=(b_x,b_y,b_z)$,
\begin{align*}
  \frac{\partial u}{\partial t}
  +u\frac{\partial u}{\partial x}
  +v\frac{\partial u}{\partial y}
  +w\frac{\partial u}{\partial z}
  &=-\frac{1}{\rho}\frac{\partial p}{\partial x}+b_x,\\[3pt]
  \frac{\partial v}{\partial t}
  +u\frac{\partial v}{\partial x}
  +v\frac{\partial v}{\partial y}
  +w\frac{\partial v}{\partial z}
  &=-\frac{1}{\rho}\frac{\partial p}{\partial y}+b_y,\\[3pt]
  \frac{\partial w}{\partial t}
  +u\frac{\partial w}{\partial x}
  +v\frac{\partial w}{\partial y}
  +w\frac{\partial w}{\partial z}
  &=-\frac{1}{\rho}\frac{\partial p}{\partial z}+b_z.
\end{align*}
\[
  \frac{\partial u}{\partial x}
  +\frac{\partial v}{\partial y}
  +\frac{\partial w}{\partial z}=0.
\]
\end{frame}

% ----- General surface stress -----
\begin{frame}{The stress vector}
At a point $\xvec$, consider a small oriented surface element with unit normal $\nvec$. The force transmitted across it is
\[
  \tvec(\nvec)\,dS,
\]
where $\tvec$ is the traction or stress vector.
\begin{itemize}\itemsep6pt
  \item In an inviscid fluid, $\tvec=-p\nvec$.
  \item In a viscous fluid, $\tvec$ may contain normal and tangential components.
  \item Reversing the normal reverses the traction: $\tvec(-\nvec)=-\tvec(\nvec)$.
\end{itemize}
\begin{punchbox}
\small Cauchy's theorem expresses the traction on every possible plane through one tensor field defined at the point.
\end{punchbox}
\end{frame}

\begin{frame}{Components of the Cauchy stress}
Define
\[
  \sigma_{ij}
  =\text{the $i$th force component on a plane whose normal is $\bm e_j$}.
\]
In Cartesian coordinates,
\[
  \stress=
  \begin{pmatrix}
  \sigma_{xx}&\sigma_{xy}&\sigma_{xz}\\
  \sigma_{yx}&\sigma_{yy}&\sigma_{yz}\\
  \sigma_{zx}&\sigma_{zy}&\sigma_{zz}
  \end{pmatrix}.
\]
Each column corresponds to a coordinate-plane normal; each row gives a force component.

\medskip
The tensor must determine the traction on an arbitrarily oriented plane.
\end{frame}

\begin{frame}{Stress on an arbitrary surface}
\begin{columns}[T,onlytextwidth]
\begin{column}{0.48\textwidth}
Applying linear momentum balance to a shrinking tetrahedron gives Cauchy's stress principle:
\[
  \boxed{\tvec(\nvec)=\stress\nvec},
  \qquad
  t_i=\sigma_{ij}n_j.
\]
Surface forces scale with length squared while inertia and body forces scale with length cubed. The surface-force balance therefore determines the limiting relation.
\end{column}
\begin{column}{0.52\textwidth}
\centering
\includegraphics[width=\linewidth]{Stress.png}

{\scriptsize Stress components and the Cauchy tetrahedron. Acheson writes the stress tensor $T_{ij}$ and the traction $\bm t$; we write $\sigma_{ij}$ and $\bm t$. Source: D. J. Acheson, \emph{Elementary Fluid Dynamics} (Oxford Univ. Press, 1990), Figs.~6.2--6.3.}
\end{column}
\end{columns}
\end{frame}

\begin{frame}{Linear momentum balance with stresses}
For a material region $\Omega(t)$,
\[
  \frac{d}{dt}\int_{\Omega(t)}\rho\uvec\,dV
  =\underbrace{\oint_{\partial\Omega(t)}\stress\nvec\,dS}_{\text{surface force}}
  +\underbrace{\int_{\Omega(t)}\rho\bvec\,dV}_{\text{body force}}.
\]
The divergence theorem gives
\[
  \oint_{\partial\Omega(t)}\stress\nvec\,dS
  =\int_{\Omega(t)}\nabla\cdot\stress\,dV.
\]
Using mass conservation (Reynolds transport theorem plus continuity),
\[
  \frac{d}{dt}\int_{V(t)}\rho\uvec\,dV
  =\int_{V(t)}\rho\frac{D\uvec}{Dt}\,dV,
\]
so that
\[
  \int_{\Omega(t)}\rho\frac{D\uvec}{Dt}\,dV
  =\int_{\Omega(t)}(\nabla\cdot\stress+\rho\bvec)\,dV.
\]
\end{frame}

\begin{frame}{Cauchy momentum equation}
Because the material region is arbitrary,
\[
  \boxed{\rho\frac{D\uvec}{Dt}
  =\nabla\cdot\stress+\rho\bvec}.
\]
For an inviscid fluid, $\stress=-p\Identity$, so
\[
  \nabla\cdot\stress=-\nabla p
\]
and Cauchy's equation reduces to the Euler momentum equation.

\begin{punchbox}
\small This is Newton's second law per unit volume. The balance law is general; a constitutive relation for $\stress$ specifies the fluid model.
\end{punchbox}
\end{frame}

\begin{frame}{Traction and momentum in index notation}
The traction relation and momentum equation use the same divergence convention:
\[
  t_i=\sigma_{ij}n_j,
  \qquad
  (\nabla\cdot\stress)_i
  =\frac{\partial\sigma_{ij}}{\partial x_j}.
\]
Hence
\[
  \boxed{
  \rho\left(
  \frac{\partial u_i}{\partial t}
  +u_j\frac{\partial u_i}{\partial x_j}
  \right)
  =\frac{\partial\sigma_{ij}}{\partial x_j}
  +\rho b_i}.
\]
\begin{punchbox}
\small The first stress index gives force direction; the second identifies the surface-normal direction.
\end{punchbox}
\end{frame}

\begin{frame}{Pressure and viscous parts of stress}
Split the Cauchy stress as
\[
  \boxed{\stress=-p\Identity+\taustress},
\]
where $-p\Identity$ is isotropic and $\taustress$ is the viscous or extra stress. Then
\[
  \rho\frac{D\uvec}{Dt}
  =-\nabla p+\nabla\cdot\taustress+\rho\bvec.
\]
For a classical fluid without couple stresses, angular-momentum balance gives
\[
  \stress=\stress^{\mathsf T}.
\]
\begin{punchbox}
\small A constitutive law for $\taustress$ closes the momentum equation.
\end{punchbox}
\end{frame}

% ----- 13. Velocity gradient -----
\begin{frame}{Deformation and rotation rates}
Decompose the velocity gradient into symmetric and antisymmetric parts:
\[
  \nabla\uvec=\Dtensor+\Wtensor,
\]
\[
  \Dtensor=\frac12\left(\nabla\uvec+(\nabla\uvec)^{\mathsf T}\right),
  \qquad
  \Wtensor=\frac12\left(\nabla\uvec-(\nabla\uvec)^{\mathsf T}\right).
\]
\begin{columns}[T,onlytextwidth]
\begin{column}{0.52\textwidth}
\hi{$\Dtensor$}: rate of stretching and shearing,
\[
  \operatorname{tr}\Dtensor=\nabla\cdot\uvec.
\]
\end{column}
\begin{column}{0.48\textwidth}
\hi{$\Wtensor$}: local rigid-body rotation rate.

For an isotropic Newtonian simple fluid, $\Wtensor$ does not contribute to viscous stress.
\end{column}
\end{columns}
\end{frame}

\begin{frame}{Geometric interpretation of the velocity gradient}
For two nearby parcels separated by $d\xvec$, their relative velocity is
\[
  d\uvec=(\nabla\uvec)\,d\xvec
  =\Dtensor\,d\xvec+\Wtensor\,d\xvec.
\]
\begin{columns}[T,onlytextwidth]
\begin{column}{0.50\textwidth}
\textbf{Symmetric part $\Dtensor$}
\begin{itemize}\itemsep5pt
  \item changes lengths;
  \item changes angles;
  \item measures stretching and shearing.
\end{itemize}
\end{column}
\begin{column}{0.50\textwidth}
\textbf{Antisymmetric part $\Wtensor$}
\begin{itemize}\itemsep5pt
  \item rotates line elements;
  \item preserves lengths and angles;
  \item represents local rigid-body spin.
\end{itemize}
\end{column}
\end{columns}
\begin{punchbox}
\small An objective viscous constitutive law cannot generate stress from a superposed rigid-body rotation alone.
\end{punchbox}
\end{frame}

\begin{frame}{Requirements for a Newtonian constitutive law}
The extra stress $\taustress$ for a Newtonian simple fluid is modeled as:
\begin{enumerate}\itemsep6pt
  \item \textbf{linear} in the local deformation rate;
  \item \textbf{local} in space and time;
  \item \textbf{objective}, so rigid-body motion alone produces no viscous stress;
  \item \textbf{isotropic}, so the material has no preferred direction.
\end{enumerate}
For an isotropic material, the available linear tensor ingredients are $\Dtensor$, its trace $\nabla\cdot\uvec$, and $\Identity$. Therefore
\[
  \taustress=2\mu\Dtensor
  +\lambda(\nabla\cdot\uvec)\Identity.
\]
\end{frame}

% ----- 14. Constitutive law -----
\begin{frame}{Newtonian constitutive law}
For a linear, local, isotropic Newtonian fluid,
\[
  \boxed{
  \taustress
  =2\mu\Dtensor
  +\lambda(\nabla\cdot\uvec)\Identity
  }.
\]
\begin{itemize}\itemsep6pt
  \item $\mu$ is the dynamic shear viscosity.
  \item $\lambda$ is the second viscosity coefficient.
  \item In three dimensions, the bulk viscosity is
  \[
    \zeta=\lambda+\frac{2}{3}\mu.
  \]
\end{itemize}
\begin{punchbox}
\small Isotropy means the response has no fixed material direction and transforms consistently under rotations.
\end{punchbox}
\end{frame}

% ----- 15. Incompressible stress -----
\begin{frame}{Incompressible Newtonian stress}
When $\nabla\cdot\uvec=0$,
\[
  \boxed{\taustress=2\mu\Dtensor},
  \qquad
  \boxed{\stress=-p\Identity+2\mu\Dtensor}.
\]
In components,
\[
  \sigma_{ij}
  =-p\delta_{ij}
  +\mu\left(
  \frac{\partial u_i}{\partial x_j}
  +\frac{\partial u_j}{\partial x_i}
  \right).
\]

\begin{columns}[T,onlytextwidth]
\begin{column}{0.50\textwidth}
Pressure supplies isotropic normal stress.
\end{column}
\begin{column}{0.50\textwidth}
Viscosity supplies shear and extensional normal stresses proportional to deformation rate.
\end{column}
\end{columns}
\end{frame}

% ----- 16. Divergence of tau -----
\begin{frame}{Viscous force density}
For constant viscosity,
\[
  \nabla\cdot(2\mu\Dtensor)
  =\mu\nabla^2\uvec
  +\mu\nabla(\nabla\cdot\uvec).
\]
Therefore, for incompressible flow,
\[
  \boxed{\nabla\cdot\taustress=\mu\nabla^2\uvec}.
\]
The Laplacian acts componentwise and diffuses momentum from regions of high velocity toward regions of low velocity.

\begin{punchbox}
\small If $\mu$ varies in space, keep $\nabla\cdot(2\mu\Dtensor)$ in conservative form; gradients of viscosity generate additional terms.
\end{punchbox}
\end{frame}

% ----- 17. NS -----
\begin{frame}{Incompressible Navier--Stokes equations}
Combining momentum balance with the incompressible Newtonian constitutive law gives
\[
\boxed{
  \rho\left[
  \frac{\partial\uvec}{\partial t}
  +(\uvec\cdot\nabla)\uvec
  \right]
  =-\nabla p+\mu\nabla^2\uvec+\rho\bvec
}
\]
\[
\boxed{\nabla\cdot\uvec=0}.
\]
Equivalently, using kinematic viscosity $\nu=\mu/\rho$,
\[
  \frac{\partial\uvec}{\partial t}
  +(\uvec\cdot\nabla)\uvec
  =-\frac{1}{\rho}\nabla p
  +\nu\nabla^2\uvec+\bvec.
\]
\end{frame}

\begin{frame}{Navier--Stokes in Cartesian partial derivatives}
\scriptsize
For $\uvec=(u,v,w)$, constant $\rho$ and $\mu$, and $\bvec=(b_x,b_y,b_z)$,
\begin{align*}
\rho\left(u_t+u u_x+v u_y+w u_z\right)
&=-p_x+\mu(u_{xx}+u_{yy}+u_{zz})+\rho b_x,\\[4pt]
\rho\left(v_t+u v_x+v v_y+w v_z\right)
&=-p_y+\mu(v_{xx}+v_{yy}+v_{zz})+\rho b_y,\\[4pt]
\rho\left(w_t+u w_x+v w_y+w w_z\right)
&=-p_z+\mu(w_{xx}+w_{yy}+w_{zz})+\rho b_z.
\end{align*}
The incompressibility constraint is
\[
  u_x+v_y+w_z=0.
\]
\begin{punchbox}
\small These are three momentum equations coupled through pressure, incompressibility, and the advective terms.
\end{punchbox}
\end{frame}

% ----- 18. Terms -----
\begin{frame}{Terms in the momentum equation}
Every term below has units of force per volume, $\mathrm{N\,m^{-3}}$.
\begin{center}
\renewcommand{\arraystretch}{1.35}
\begin{tabular}{ll}
\toprule
Term & Physical role \\
\midrule
$\displaystyle \rho\frac{\partial\uvec}{\partial t}$
& local inertia \\
$\displaystyle \rho(\uvec\cdot\nabla)\uvec$
& advective inertia \\
$\displaystyle -\nabla p$
& pressure-gradient force \\
$\displaystyle \mu\nabla^2\uvec$
& viscous momentum diffusion \\
$\displaystyle \rho\bvec$
& body force \\
\bottomrule
\end{tabular}
\end{center}
\begin{punchbox}
\centering
\small The Reynolds number compares the advective-inertia scale $\rho U^2/L$ with the viscous scale $\mu U/L^2$.
\end{punchbox}
\end{frame}

% ----- 19. Index form -----
\begin{frame}{Cartesian index form}
Using Einstein summation (every repeated spatial index is summed from 1 to 3),
\[
  \boxed{
  \rho\left(
  \frac{\partial u_i}{\partial t}
  +u_j\frac{\partial u_i}{\partial x_j}
  \right)
  =-\frac{\partial p}{\partial x_i}
  +\mu\frac{\partial^2u_i}{\partial x_j\partial x_j}
  +\rho b_i
  },
  \qquad i=1,2,3,
\]
with
\[
  \boxed{\frac{\partial u_i}{\partial x_i}=0}.
\]

\begin{punchbox}
\small Index notation keeps the vector structure visible without placing three long component equations on one slide.
\end{punchbox}
\end{frame}

\begin{frame}{Steady incompressible Navier--Stokes equations}
For a steady flow, $\partial\uvec/\partial t=0$. The equations become
\[
  \boxed{
  \rho(\uvec\cdot\nabla)\uvec
  =-\nabla p+\mu\nabla^2\uvec+\rho\bvec},
  \qquad
  \boxed{\nabla\cdot\uvec=0}.
\]
\begin{itemize}\itemsep6pt
  \item Parcel acceleration may remain nonzero through advection.
  \item The equations remain nonlinear unless the assumed geometry eliminates $(\uvec\cdot\nabla)\uvec$.
  \item Fully developed channel and pipe flows provide useful analytical examples.
\end{itemize}
\end{frame}

% ----- 20. Parallel flow -----
\begin{frame}{Fully developed unidirectional flow}
Assume a parallel flow
\[
  \uvec=(u(y,t),0,0),
  \qquad
  \frac{\partial u}{\partial x}
  =\frac{\partial u}{\partial z}=0.
\]
Incompressibility is automatic, and the advective term vanishes. The streamwise equation reduces to
\[
  \boxed{
  \rho\frac{\partial u}{\partial t}
  =-\frac{\partial p}{\partial x}
  +\mu\frac{\partial^2u}{\partial y^2}
  +\rho b_x
  }.
\]

\begin{columns}[T,onlytextwidth]
\begin{column}{0.50\textwidth}
\hi{Couette flow}: wall motion drives $u(y)$.
\end{column}
\begin{column}{0.50\textwidth}
\hi{Poiseuille flow}: a pressure gradient drives $u(y)$.
\end{column}
\end{columns}
\end{frame}

\begin{frame}{Two-dimensional incompressible flow}
Let $\uvec=(u(x,y,t),v(x,y,t),0)$. Then
\[
  \rho(u_t+u u_x+v u_y)
  =-p_x+\mu(u_{xx}+u_{yy})+\rho b_x,
\]
\[
  \rho(v_t+u v_x+v v_y)
  =-p_y+\mu(v_{xx}+v_{yy})+\rho b_y,
\]
with
\[
  \boxed{u_x+v_y=0}.
\]
\begin{punchbox}
\small Two-dimensional models support nonlinear advection, separation, and vortex shedding, but omit three-dimensional instabilities and vortex stretching.
\end{punchbox}
\end{frame}

% ----- Boundary conditions -----
\begin{frame}{Boundary conditions for Navier--Stokes}
The equations must be paired with conditions on $\partial\Omega$ that represent the physical boundary.
\begin{itemize}\itemsep7pt
  \item A solid wall constrains fluid motion relative to the wall.
  \item An inflow specifies incoming velocity or flux.
  \item An open boundary controls traction or another outflow model.
  \item An interface couples kinematics to a traction balance.
  \item Periodic faces identify values across repeating boundaries.
\end{itemize}
\begin{punchbox}
\small Boundary conditions are part of the physical and mathematical model. Changing them can change the solution even when the differential equations remain the same.
\end{punchbox}
\end{frame}

% ----- 21. Data -----
\begin{frame}{Initial and boundary data}
An unsteady incompressible problem needs an initial velocity
\[
  \uvec(\xvec,0)=\uvec_0(\xvec),
  \qquad
  \nabla\cdot\uvec_0=0,
\]
compatible with the boundary data at $t=0$.

Partition the boundary into portions with velocity and traction data:
\[
  \partial\Omega=\Gamma_D\cup\Gamma_N,
  \qquad
  \Gamma_D\cap\Gamma_N=\varnothing,
\]
\[
  \uvec=\overline{\uvec}
  \ \text{on }\Gamma_D,
  \qquad
  \stress\nvec=\overline{\tvec}
  \ \text{on }\Gamma_N.
\]

For velocity prescribed on the entire boundary, incompressibility requires
\[
  \int_{\partial\Omega}\overline{\uvec}\cdot\nvec\,dS=0.
\]
\end{frame}

% ----- 22. Walls and inflow -----
\begin{frame}{Solid-wall and inflow conditions}
\begin{columns}[T,onlytextwidth]
\begin{column}{0.50\textwidth}
\textbf{No-slip solid wall}
\[
  \uvec=\mathbf{U}_{\mathrm{wall}}.
\]
For a stationary wall, $\uvec=\mathbf{0}$.

\medskip
\textbf{No penetration only}
\[
  (\uvec-\mathbf{U}_{\mathrm{wall}})\cdot\nvec=0.
\]
This constrains the normal velocity component but still requires a tangential condition.
\end{column}
\begin{column}{0.50\textwidth}
\textbf{Prescribed inflow}
\[
  \uvec=\uvec_{\mathrm{in}}(\xvec,t).
\]
The profile should be physically plausible and compatible with the net outflow.

\medskip
\begin{punchbox}
\small Dirichlet data specify velocity. No penetration is an essential condition on one component, not a traction condition.
\end{punchbox}
\end{column}
\end{columns}
\end{frame}

% ----- 23. Traction -----
\begin{frame}{Traction and open-boundary conditions}
A traction condition prescribes force per unit area:
\[
  \stress\nvec=\overline{\tvec}.
\]
Examples include:
\begin{itemize}\itemsep6pt
  \item stress-free idealization: $\stress\nvec=\mathbf{0}$;
  \item known external pressure with negligible external shear:
  \[
    \stress\nvec=-p_{\mathrm{ext}}\nvec;
  \]
  \item a variational ``do-nothing'' outflow, whose precise traction depends on the chosen weak form.
\end{itemize}

\begin{punchbox}
\small A zero normal velocity gradient is not a universal outflow condition. Open boundaries require a condition consistent with the flow direction, stress model, and numerical formulation.
\end{punchbox}
\end{frame}

\begin{frame}{Robin and Navier-slip conditions}
A Robin condition combines a velocity response with a traction response. A common example is Navier slip:
\[
  (\uvec-\mathbf U_{\mathrm{wall}})\cdot\nvec=0,
\]
\[
  (\stress\nvec)_{\parallel}
  +\frac{\mu}{\ell_s}
  (\uvec-\mathbf U_{\mathrm{wall}})_{\parallel}=0.
\]
Here $\ell_s$ is the slip length and $\nvec$ is the outward normal of the fluid domain.

\begin{itemize}\itemsep5pt
  \item $\ell_s\to0$ approaches no slip.
  \item Large $\ell_s$ approaches small tangential traction.
  \item Resistive or porous boundaries use related mixed relations.
\end{itemize}
\end{frame}

% ----- 24. Other BCs -----
\begin{frame}{Slip, periodic, and free-surface conditions}
\begin{columns}[T,onlytextwidth]
\begin{column}{0.52\textwidth}
\textbf{Navier slip}
\[
  (\uvec-\mathbf{U}_{\mathrm{wall}})\cdot\nvec=0,
\]
\[
  (\stress\nvec)_{\parallel}
  +\frac{\mu}{\ell_s}
  (\uvec-\mathbf{U}_{\mathrm{wall}})_{\parallel}=0.
\]
The displayed sign uses the outward normal of the fluid domain.

\medskip
\textbf{Periodic flow}

Velocity is periodic. Pressure may be periodic up to a prescribed pressure drop.
\end{column}
\begin{column}{0.48\textwidth}
\textbf{Free surface}

Two conditions are needed:
\begin{itemize}\itemsep5pt
  \item a kinematic condition that moves the interface with the fluid;
  \item a dynamic condition that balances the traction jump with surface-tension and Marangoni forces.
\end{itemize}
\begin{punchbox}
\small State the normal and curvature conventions when writing the surface-tension sign.
\end{punchbox}
\end{column}
\end{columns}
\end{frame}

\begin{frame}{Choosing boundary conditions}
Match each condition to the physics and to the mathematical formulation:
\begin{itemize}\itemsep6pt
  \item use velocity data where motion is known, such as an inflow or no-slip wall;
  \item use traction or a formulation-specific open condition where force exchange is known;
  \item use periodic data only when the geometry and forcing repeat;
  \item state both kinematic and dynamic conditions at a free surface;
  \item test outflow placement and sensitivity because artificial open boundaries are delicate.
\end{itemize}
\begin{punchbox}
\small Do not prescribe more information than the coupled velocity-pressure system can accept. Check flux compatibility and the pressure gauge.
\end{punchbox}
\end{frame}

% ----- 25. Pressure -----
\begin{frame}{Pressure level and pressure boundary data}
In the momentum equation, pressure appears through $\nabla p$.
\begin{itemize}\itemsep6pt
  \item With velocity prescribed everywhere, $p$ is determined only up to an additive constant. Fix a gauge such as
  \[
    \int_\Omega p\,dV=0
    \qquad\text{or}\qquad
    p(\xvec_0)=0.
  \]
  \item A traction boundary can determine the pressure level because $p$ is part of $\stress\nvec$.
  \item A ``pressure outlet'' is normally shorthand for a traction model; viscous and tangential stresses must be handled consistently.
\end{itemize}

\begin{punchbox}
\small Prescribing a complete velocity and a complete traction on the same boundary segment generally overdetermines the standard incompressible formulation.
\end{punchbox}
\end{frame}

% ----- 26. Pressure Poisson -----
\begin{frame}{Pressure Poisson equation}
Take the divergence of the constant-density, constant-viscosity momentum equation and use $\nabla\cdot\uvec=0$:
\[
  \boxed{
  \nabla^2p
  =-\rho\,\nabla\cdot\left[(\uvec\cdot\nabla)\uvec\right]
  +\rho\,\nabla\cdot\bvec
  }.
\]
For constant gravity, $\nabla\cdot\bvec=0$.

The pressure equation inherits boundary information from the normal momentum balance and the chosen velocity or traction conditions.

\begin{punchbox}
\small Pressure acts as a Lagrange multiplier enforcing incompressibility. The Poisson equation is a consequence of momentum balance, not an independent constitutive law.
\end{punchbox}
\end{frame}

% ----- 27. Checklist -----
\begin{frame}{Boundary-data checklist}
\begin{enumerate}\itemsep7pt
  \item Match the condition to the physical boundary: wall, inflow, outflow, interface, or periodic face.
  \item Avoid prescribing full velocity and full traction on the same boundary segment.
  \item Enforce incompressible flux compatibility for prescribed velocity data.
  \item Supply divergence-free, boundary-compatible initial data for unsteady flow.
  \item Fix the pressure gauge only when the boundary conditions leave it undetermined.
\end{enumerate}
\begin{punchbox}
\centering
\small Boundary conditions are part of the mathematical model; small changes can produce different physical solutions.
\end{punchbox}
\end{frame}

\begin{frame}{Well-posed and overdetermined examples}
\small
\textbf{Compatible examples}
\begin{itemize}\itemsep4pt
  \item prescribed inflow velocity with a consistent traction condition at outflow;
  \item no-slip walls with traction specified on the remaining open boundary;
  \item periodic velocity with a pressure gauge or a prescribed mean pressure drop.
\end{itemize}

\medskip
\textbf{Common problems}
\begin{itemize}\itemsep4pt
  \item prescribing complete velocity and complete traction on the same boundary segment;
  \item prescribing incompatible net inflow when velocity is fixed everywhere;
  \item treating a scalar pressure value as a complete open-boundary condition without specifying the accompanying viscous traction model.
\end{itemize}
\begin{punchbox}
\small Well-posedness also depends on domain geometry, regularity, and the exact weak or numerical formulation; these examples are modeling guides, not a general theorem.
\end{punchbox}
\end{frame}

% ----- 28. Scope -----
\begin{frame}{Assumptions and extensions}
The derived Navier--Stokes form assumes:
\begin{itemize}\itemsep5pt
  \item a continuum fluid;
  \item incompressibility and uniform constant density;
  \item a Newtonian, isotropic constitutive law;
  \item constant dynamic viscosity;
  \item classical Cauchy stress without couple stresses.
\end{itemize}

Extensions modify specific parts of the derivation:
\begin{itemize}\itemsep4pt
  \item compressible flow couples momentum to density and energy;
  \item variable or non-Newtonian viscosity changes $\nabla\cdot\taustress$;
  \item moving interfaces add kinematic and traction-jump conditions.
\end{itemize}
\end{frame}

% ----- 29. Summary -----
\begin{frame}{Derivation summary}
\begin{enumerate}\itemsep8pt
  \item Parcel acceleration is $D\uvec/Dt=\partial_t\uvec+(\uvec\cdot\nabla)\uvec$.
  \item Mass conservation gives continuity; incompressibility gives $\nabla\cdot\uvec=0$.
  \item Surface force is represented by traction $\tvec=\stress\nvec$.
  \item Momentum balance gives $\rho D\uvec/Dt=\nabla\cdot\stress+\rho\bvec$.
  \item The incompressible Newtonian law $\stress=-p\Identity+2\mu\Dtensor$ yields Navier--Stokes.
\end{enumerate}
\begin{punchbox}
\centering
\small Conservation laws determine the balances. The constitutive law and boundary data determine the particular fluid model and problem.
\end{punchbox}
\end{frame}

% ----- 30. References -----
\begin{frame}{References}
\small
\begin{itemize}\itemsep7pt
  \item D. J. Acheson, \emph{Elementary Fluid Dynamics}. Clarendon Press, 1990.
  \item G. K. Batchelor, \emph{An Introduction to Fluid Dynamics}. Cambridge University Press, 1967.
  \item L. D. Landau and E. M. Lifshitz, \emph{Fluid Mechanics}, 2nd ed. Pergamon Press, 1987.
  \item M. E. Gurtin, E. Fried, and L. Anand, \emph{The Mechanics and Thermodynamics of Continua}. Cambridge University Press, 2010.
  \item C. S. Peskin, ``The immersed boundary method,'' \emph{Acta Numerica} 11 (2002), \href{https://doi.org/10.1017/S0962492902000077}{doi:10.1017/S0962492902000077}.
\end{itemize}
\end{frame}

% ----- 31. Further study -----
\begin{frame}{Further study}
\small
\begin{itemize}\itemsep7pt
  \item \href{https://ocw.mit.edu/courses/1-060-engineering-mechanics-ii-spring-2006/}{MIT OpenCourseWare 1.060: Engineering Mechanics II} for an introductory treatment.
  \item \href{https://ocw.mit.edu/courses/2-25-advanced-fluid-mechanics-fall-2013/}{MIT OpenCourseWare 2.25: Advanced Fluid Mechanics} for graduate-level notes and problems.
  \item \href{https://sciencedemonstrations.fas.harvard.edu/fluid-dynamics}{Harvard Natural Sciences Lecture Demonstrations: Fluid Dynamics} for physical demonstrations.
  \item \href{https://onlinecourses.nptel.ac.in/noc22_ge11/preview}{NPTEL: Fundamentals of Fluid Mechanics for Chemical and Biomedical Engineers} for biomedical examples.
  \item \href{https://www.flowvis.org/}{University of Colorado Boulder Flow Visualization} for experimental and artistic flow imagery.
\end{itemize}
\end{frame}

% ----- 32. Appendix: operators -----
\begin{frame}{Appendix: Cartesian vector operators}
For $\uvec=(u,v,w)$ and a scalar $\phi$,
\begin{align*}
  \nabla\phi
  &=\left(
  \frac{\partial\phi}{\partial x},
  \frac{\partial\phi}{\partial y},
  \frac{\partial\phi}{\partial z}
  \right),\\[5pt]
  \nabla\cdot\uvec
  &=\frac{\partial u}{\partial x}
  +\frac{\partial v}{\partial y}
  +\frac{\partial w}{\partial z},\\[5pt]
  \nabla\times\uvec
  &=\left(
  \frac{\partial w}{\partial y}-\frac{\partial v}{\partial z},
  \frac{\partial u}{\partial z}-\frac{\partial w}{\partial x},
  \frac{\partial v}{\partial x}-\frac{\partial u}{\partial y}
  \right),\\[5pt]
  \nabla^2\uvec
  &=\left(\nabla^2u,\nabla^2v,\nabla^2w\right).
\end{align*}
The vorticity is $\bm{\omega}=\nabla\times\uvec$.
\end{frame}

% ----- 33. Appendix: identities -----
\begin{frame}{Appendix: stress-divergence identities}
With the traction convention $\tvec=\stress\nvec$,
\[
  (\nabla\cdot\stress)_i
  =\frac{\partial\sigma_{ij}}{\partial x_j},
  \qquad
  \oint_{\partial\Omega}\stress\nvec\,dS
  =\int_\Omega\nabla\cdot\stress\,dV.
\]
For constant $\mu$,
\[
  \nabla\cdot(2\mu\Dtensor)
  =\mu\nabla^2\uvec
  +\mu\nabla(\nabla\cdot\uvec).
\]
For the full compressible Newtonian law with constant $\lambda$ and $\mu$,
\[
  \nabla\cdot\taustress
  =\mu\nabla^2\uvec
  +(\lambda+\mu)\nabla(\nabla\cdot\uvec).
\]
\end{frame}

\end{document}
